\exercisesection{INTEGRAL CLOSURE OF A COMPLETE LOCAL RING}

\begin{enumerate}
\def\labelenumi{\arabic{enumi})}
\tightlist
\item
  Let A be a semi-local Noetherian integral domain of dimension 1.
\end{enumerate}

\begin{enumerate}
\def\labelenumi{\alph{enumi})}
\item
  The integral closure of A is a finite A-algebra if and only if the completion \(\hat{\mathbf{A}}\) of A is reduced. (If \(\hat{\mathbf{A}}\) is reduced, use the corollary of Theorem 2 of § 4, No. 2. For the converse implication, reduce to the case where A is integrally closed, and argue as in the proof of Theorem 3 of § 4, No. 4).
\item
  The following three conditions are equivalent:
\end{enumerate}

\begin{enumerate}
\def\labelenumi{(\roman{enumi})}
\item
  A is a Japanese ring;
\item
  A is a Nagata ring;
\item
  \(\hat{\mathbf{A}}\) is reduced and if \(q_{1},\ldots,q_{n}\) are the minimal prime ideals of \(\hat{\mathbf{A}}\) then the fields \(\kappa(q_{i})\) are separable extensions of the field of fractions of A.
\end{enumerate}

\begin{enumerate}
\def\labelenumi{\arabic{enumi})}
\setcounter{enumi}{1}
\tightlist
\item
  Let A be a Noetherian local ring of dimension 1. Then the following conditions are equivalent:
\end{enumerate}

\begin{enumerate}
\def\labelenumi{\alph{enumi})}
\item
  A is regular;
\item
  A is integrally closed;
\item
  A is a discrete valuation ring.
\end{enumerate}

\begin{enumerate}
\def\labelenumi{\arabic{enumi})}
\setcounter{enumi}{2}
\tightlist
\item
  A ring A is said to be normal if the ring \(\mathbf{A}_{\mathfrak{p}}\) is integrally closed for every prime ideal p of A.
\end{enumerate}

\begin{enumerate}
\def\labelenumi{\alph{enumi})}
\item
  A normal ring A is reduced, and every prime ideal p of A contains exactly one minimal prime ideal of A. For every minimal prime ideal p of A, the ring A/p is integrally closed.
\item
  A Noetherian normal ring A is isomorphic to a finite product of integrally closed Noetherian rings.
\item
  Let A be a Noetherian normal ring and let a be an element of A. Then Ass \(_{A}\) (A/aA) contains no embedded prime ideal (IV, § 2, no. 3, remark).
\end{enumerate}

\begin{enumerate}
\def\labelenumi{\arabic{enumi})}
\setcounter{enumi}{3}
\tightlist
\item
  Let A be a ring. We consider the following two conditions on A:
\end{enumerate}

(R1) For every prime ideal \(\mathfrak{p}\) of A of height ≤ 1, the ring \(A_{\mathfrak{p}}\) is regular.

(S2) The set \(\mathrm{Ass}_{\mathbf{A}}(\mathbf{A})\) and, for every cancellable element \(a\) of A, the set \(\mathrm{Ass}_{\mathbf{A}}(\mathbf{A}/a\mathbf{A})\) do not contain an embedded prime ideal.

A normal Noetherian ring (exercise 3) satisfies (R1) and (S2). (Use exercises 2 and 3, c.)

\begin{enumerate}
\def\labelenumi{\arabic{enumi})}
\setcounter{enumi}{4}
\tightlist
\item
  Let A be a Noetherian ring satisfying conditions (R1) and (S2).
\end{enumerate}

\begin{enumerate}
\def\labelenumi{\alph{enumi})}
\item
  Show that A is reduced.
\item
  Let \(a\) be a cancellable element of A. If an element \(x\) of A is such that its image in \(\mathbf{A}_{\mathfrak{p}}\) belongs to \(a\mathbf{A}_{\mathfrak{p}}\) for every \(\mathfrak{p} \in \operatorname{Ass}_{\mathbf{A}}(\mathbf{A}/a\mathbf{A})\) , then \(x\) belongs to \(a\mathbf{A}\) . (Use IV, § 2, no. 3, prop. 5.)
\item
  Prove that A is integrally closed in its total ring of fractions R. (If \(a, b, c_{1}, \ldots, c_{n}\) are elements of A satisfying
\end{enumerate}

\begin{enumerate}
\def\labelenumi{(\roman{enumi})}
\item
  \(b\) is cancellable in A,
\item
  \((b^{-1}a)^{n} + c_{1}(b^{-1}a)^{n-1} + \cdots + c_{n} = 0\) in R, then prove successively that we have \(a \in bA_{\mathfrak{p}}\) for every prime ideal \(\mathfrak{p}\) of height 1 in A, and then that \(a \in bA\) .)
\end{enumerate}

\begin{enumerate}
\def\labelenumi{\alph{enumi})}
\setcounter{enumi}{3}
\tightlist
\item
  Prove that A is normal. (Note that an idempotent of R is integral over A.)
\end{enumerate}

\begin{enumerate}
\def\labelenumi{\arabic{enumi})}
\setcounter{enumi}{5}
\tightlist
\item
  \begin{enumerate}
  \def\labelenumii{\alph{enumii})}
  \tightlist
  \item
    Let A be a ring, and let M be a faithful Noetherian A-module. Show that A is a Noetherian ring.
  \end{enumerate}
\end{enumerate}

\begin{enumerate}
\def\labelenumi{\alph{enumi})}
\setcounter{enumi}{1}
\item
  Let A be a ring and M a faithful finitely generated A-module. Suppose that for every increasing sequence of ideals \((\mathfrak{a}_{i})_{i\in\mathbb{N}}\) of A, the sequence of submodules \((\mathfrak{a}_{i}\mathbf{M})_{i\in\mathbb{N}}\) of M is stationary. Then A is a Noetherian ring. (By (a), it suffices to prove that M is a Noetherian A-module. Reduce to the case where, for every nonzero ideal a of A, the A-module M/aM is Noetherian, and where, for every nonzero submodule N of M, the A-module M/N is not faithful.)
\item
  Let B be a Noetherian ring and A a subring of B such that B is a finite A-algebra. Then A is a Noetherian ring. (This allows one to remove the hypothesis that A be Noetherian in prop. 2 of § 4, no. 1.)
\end{enumerate}

\begin{enumerate}
\def\labelenumi{\arabic{enumi})}
\setcounter{enumi}{6}
\tightlist
\item
  Let A be a Krull ring and let P be the set of its height-1 prime ideals. Suppose that for every prime ideal \(p \in P\) the ring A/p is Noetherian, and we denote by K the field of fractions of A. Prove that A is Noetherian.
\end{enumerate}

(Let \(\mathfrak{p} \in \mathrm{P}\) and \(x \in \mathbf{K}\) such that \(v_{\mathfrak{p}}(x) = 1\) and \(v_{\mathfrak{q}}(x) \leqslant 0\) for \(\mathfrak{q} \in \mathrm{P} - \{\mathfrak{p}\}\) . Put \(\mathbf{B} = \mathbf{A}[x]\) . Then we have the following properties:

\begin{enumerate}
\def\labelenumi{(\roman{enumi})}
\item
  \(\mathfrak{p} = x\mathbf{B}\cap \mathbf{A};\)
\item
  the inclusion of A into B induces, by passing to quotients, an isomorphism from A/p onto B/xB;
\item
  for every integer \(n \geqslant 0\) , the ring \(\mathbf{B}/x^{n}\mathbf{B}\) is Noetherian;
\item
  for every integer \(n \geqslant 0\) , the ring \(\mathbf{A}/(x^n\mathbf{B} \cap \mathbf{A})\) is Noetherian (use Exercise 6);
\item
  for every integer \(n \geqslant 0\) , the A-module \(\mathbf{A}/\mathfrak{p}^{(n)}\) is Noetherian. Recall that we set \(\mathfrak{p}^{(n)} = \mathbf{A} \cap \mathfrak{p}^n\mathbf{A}_{\mathfrak{p}}\) (cf. IV, § 2, exerc. 18.)
\end{enumerate}

\begin{enumerate}
\def\labelenumi{\arabic{enumi})}
\setcounter{enumi}{7}
\tightlist
\item
  Let A be an integral domain, K its field of fractions. Let B be a ring and \(\varphi: A \to B\) a ring homomorphism making B into a faithfully flat A-module. Assume that B has only finitely many minimal prime ideals \(\mathfrak{p}_1, \ldots, \mathfrak{p}_n\) . For \(1 \leqslant i \leqslant n\) , we set \(\mathbf{B}_i = \mathbf{B} / \mathfrak{p}_i\) , denote by \(\mathbf{L}_i\) the field of fractions of \(\mathbf{B}_i\) and L the product ring of \(\mathbf{L}_i\) . Finally, we denote \(\overline{\mathbf{A}}\) the integral closure of A, \(\overline{\mathbf{B}}_i\) that of \(\mathbf{B}_i\) .
\end{enumerate}

\begin{enumerate}
\def\labelenumi{\alph{enumi})}
\item
  The homomorphism from A to L deduced from \(\varphi\) is injective, and extends to an injective homomorphism \(\psi\) from K to L.
\item
  We have \(\overline{\mathbf{A}} = \psi^{-1}(\prod_{i=1}^{n} \overline{\mathbf{B}}_i)\) . (One may prove that if \(x \in \psi^{-1}(\prod_{i=1}^{n} \overline{\mathbf{B}}_i)\) and if \(x'\) is its image in \(\mathbf{K} \otimes_{\mathbf{A}} \mathbf{B}\) , there exists a monic polynomial \(\mathbf{P}\) with coefficients in \(\mathbf{B}\) such that \(\mathbf{P}(x')\) is nilpotent in \(\mathbf{K} \otimes_{\mathbf{A}} \mathbf{B}\) .)
\end{enumerate}

\begin{enumerate}
\def\labelenumi{\arabic{enumi})}
\setcounter{enumi}{8}
\tightlist
\item
  Let A be a Noetherian integral domain, K its field of fractions, E a finite extension of K, and \(\overline{\mathbf{A}}\) the integral closure of A in E.
\end{enumerate}

\begin{enumerate}
\def\labelenumi{\alph{enumi})}
\item
  If A is a local ring with completion \(\hat{\mathbf{A}}\), then \((\hat{\mathbf{A}} \otimes_{\mathbf{A}} \overline{\mathbf{A}})_{\mathrm{red}}\) is a finite \(\hat{\mathbf{A}}\)-algebra \(^{1}\) (Reduce to the case where \(\mathbf{K} = \mathbf{E}\) . Using III, § 3, no. 4, cor. 2 of th. 3, prove that the nonzero elements of A are not zero divisors in \(\hat{\mathbf{A}}_{\mathrm{red}}\) . Then prove that there exists a \(\hat{\mathbf{A}}_{\mathrm{red}}\)-injective homomorphism from \((\hat{\mathbf{A}} \otimes_{\mathbf{A}} \overline{\mathbf{A}})_{\mathrm{red}}\) into the total ring of fractions of \(\hat{\mathbf{A}}_{\mathrm{red}}\) . Conclude.)
\item
  For every prime ideal p of A, the \(\kappa(p)\) -algebra \((\overline{A} \otimes_A \kappa(p))_{\text{red}}\) is finite. (Reduce to the case where A is local, with maximal ideal p. Then observe that \((\overline{A} \otimes_A \kappa(p))_{\text{red}}\) is a quotient of \((\overline{A} \otimes_A \hat{A})_{\text{red}} \otimes \kappa(p).)\)
\end{enumerate}

\begin{enumerate}
\def\labelenumi{\arabic{enumi})}
\setcounter{enumi}{9}
\tightlist
\item
  Let A be a local Noetherian integral domain, \(\hat{\mathbf{A}}\) its completion, \(\mathfrak{p}_1, \ldots, \mathfrak{p}_n\) the minimal prime ideals of \(\hat{\mathbf{A}}\) , and for \(1 \leqslant i \leqslant n\) , set \(\mathbf{B}_i = \hat{\mathbf{A}} / \mathfrak{p}_i\) .
\end{enumerate}

\begin{enumerate}
\def\labelenumi{\alph{enumi})}
\item
  The ring \(\mathbf{B}_i\) is Japanese.
\item
  The integral closure \(\overline{\mathbf{B}}_{i}\) of \(\mathbf{B}_{i}\) is a Krull ring.
\item
  The integral closure \(\overline{\mathbf{A}}\) of \(\mathbf{A}\) is a Krull ring. (Use Exercise 8 and VII, § 1, No. 3, Examples 3 and 4.)
\item
  \(\overline{\mathrm{A}}\) has only a finite number of maximal ideals, and their residue fields are finite extensions of that of A. (Use Exercise 9, b).
\item
  For every prime ideal p of A, there are only finitely many prime ideals P of \(\overline{A}\) above p, and the fields \(\kappa(\mathfrak{P})\) are finite extensions of \(\kappa(p)\) (Reduce to the case where A is local with maximal ideal p.)
\end{enumerate}

\begin{enumerate}
\def\labelenumi{\arabic{enumi})}
\setcounter{enumi}{10}
\tightlist
\item
  Let A be a Noetherian integral domain, K its field of fractions, E a finite extension of K, \(\overline{A}\) the integral closure of A in E.
\end{enumerate}

\begin{enumerate}
\def\labelenumi{\alph{enumi})}
\item
  For every prime ideal p of A, there are only finitely many prime ideals P of \(\overline{A}\) above p and the fields \(\kappa(\mathfrak{P})\) are finite extensions of \(\kappa(p)\) . (Reduce to the case where A is local with maximal ideal p and where E = K, and use Exerc. 10.)
\item
  Let P be the set of prime ideals of height 1 of \(\overline{\mathbf{A}}\) . If \(p \in P\) , then \(\overline{\mathbf{A}}_p\) is a discrete valuation ring.
\item
  We have \(\overline{\mathbf{A}} = \bigcap_{\mathfrak{p} \in \mathbb{P}} \overline{\mathbf{A}}_{\mathfrak{p}}\) (in E).
\end{enumerate}

\begin{enumerate}
\def\labelenumi{\arabic{enumi})}
\setcounter{enumi}{11}
\tightlist
\item
  Let A be an integral domain, \(\overline{\mathbf{A}}\) its integral closure. Let \(a_1, \ldots, a_r\) be nonzero elements of \(\overline{\mathbf{A}}\) , and let b be the fractional ideal \(\sum_{i=1}^{r} \mathbf{A}a_i\) of A.
\end{enumerate}

\begin{enumerate}
\def\labelenumi{\alph{enumi})}
\item
  There exists a fractional ideal \(\mathfrak{a}\) of A such that \(\mathfrak{a}a_{i} \subset \mathfrak{a}\) for \(1 \leqslant i \leqslant r\). We also have \(\tilde{\mathfrak{a}}a_{i} \subset \tilde{\mathfrak{a}}\) for \(1 \leqslant i \leqslant r\) ; we recall (VII, § 1, no. 1) that \(\tilde{\mathfrak{a}}\) is the intersection of the fractional principal ideals containing \(\mathfrak{a}\) .
\item
  Let y be a nonzero element of \(\tilde{\mathfrak{b}}\) . Then
\end{enumerate}

\[
\mathrm{A} y ^ {- 1} \supset \bigcap_ {i = 1} ^ {r} \mathrm{A} a _ {i} ^ {- 1} \quad \text { and } \quad \widetilde {\mathfrak {a}} \subset \widetilde {\mathfrak {a}} y ^ {- 1}.
\]

\begin{enumerate}
\def\labelenumi{\alph{enumi})}
\setcounter{enumi}{2}
\item
  We have \(\tilde{\mathbf{b}}\subset \overline{\mathbf{A}}\)
\item
  Let z be an element of \(\overline{\mathbf{A}}:\overline{\mathbf{A}}\mathbf{b}\) . Then \(\overline{\mathbf{A}}\) contains \(\tilde{\mathbf{b}}z\) .
\item
  We have b ⊂ \(\overline{A}:(\overline{A}:\overline{A}b)\) .
\item
  Let \(\mathfrak{P}\) be a divisorial ideal of \(\overline{\mathbf{A}}\) and \(\mathfrak{p} = \mathfrak{P} \cap \mathbf{A}\) . Then \(\mathfrak{p}\) is divisorial in \(\mathbf{A}\) . (By \(e\) ), we obtain \(\tilde{\mathfrak{p}} \subset \mathfrak{P}\) . But we also have \(\tilde{\mathfrak{p}} \subset \mathbf{A}\) .)
\end{enumerate}

\begin{enumerate}
\def\labelenumi{\arabic{enumi})}
\setcounter{enumi}{12}
\tightlist
\item
  Let A be a Noetherian integral domain, K its field of fractions, \(\overline{\mathbf{A}}\) its integral closure, \(f\) a nonzero element of A, \(\mathfrak{P}\) a prime ideal of height 1 of \(\overline{\mathbf{A}}\) containing \(f\), and let \(\mathfrak{p} = \mathbf{A} \cap \mathfrak{P}\) be the prime ideal
\end{enumerate}

\(^{1}\) For every ring B, we denote by \(B_{red}\) the quotient of B by the ideal of nilpotent elements of B.

of A. Prove that one has \(\mathfrak{p} \in \operatorname{Ass}_{\mathbf{A}}(\mathbf{A}/f\mathbf{A})\) by reducing to the case where A is local, with maximal ideal \(\mathfrak{p}\), and by establishing successively under this hypothesis the following assertions: a) \(\mathfrak{P}\) is divisorial (use Exercise 10, c)).

\begin{enumerate}
\def\labelenumi{\alph{enumi})}
\setcounter{enumi}{1}
\item
  p is divisorial (use exerc. 12, f)).
\item
  If \(a_{1},\ldots,a_{r}\) are nonzero elements of K such that
\end{enumerate}

\[
\mathrm{A}: \mathfrak {p} = \mathrm{A} + \sum_ {i = 1} ^ {r} \mathrm{A} a _ {i}, \quad \text { and } \quad \mathfrak {p} = \bigcap_ {i = 1} ^ {r} (\mathrm{A} \cap a _ {i} ^ {- 1} \mathrm{A})
\]

and there exists an index \(i\) such that \(\mathfrak{p} = \mathbf{A} \cap a_i^{-1}\mathbf{A}\) .

\begin{enumerate}
\def\labelenumi{\alph{enumi})}
\setcounter{enumi}{3}
\item
  There exists a nonzero element g of A such that \(\mathfrak{p} \in \operatorname{Ass}_{\mathbf{A}}(\mathbf{A}/g\mathbf{A})\) .
\item
  \(\mathfrak{p} \in \operatorname{Ass}_{\mathbf{A}}(\mathbf{A}/f\mathbf{A})\) .
\end{enumerate}

\begin{enumerate}
\def\labelenumi{\arabic{enumi})}
\setcounter{enumi}{13}
\item
  Let A be a Noetherian integral domain, K its field of fractions, \(\overline{\mathbf{A}}\) its integral closure in a finite extension E of K. Then \(\overline{\mathbf{A}}\) is a Krull ring. (Use Exerc. 11, c) and, to prove that every element \(f\) of \(\overline{\mathbf{A}}\) , nonzero, belongs to only finitely many prime ideals of height 1, reduce to the case where \(f \in \mathbf{A}\) , \(\mathbf{E} = \mathbf{K}\) and use Exercises 13, b) and 11, a).
\item
  Let A be a Noetherian ring, B an integral A-algebra having only finitely many minimal prime ideals \(p_{1}, \ldots, p_{n}\) . For \(1 \leqslant i \leqslant n\) , denote by \(q_{i}\) the inverse image of \(p_{i}\) in A, and suppose that \(\kappa(p_{i})\) is a finite extension of \(\kappa(q_{i})\) . Then, for every element x of B, the topological subspace V(x) of Spec(B) has only finitely many irreducible components. (Reduce to the case where B is integrally closed and contains A, and use exerc. 14.)
\item
  Let A be a Noetherian integral domain, \(\overline{\mathbf{A}}\) its integral closure in a finite extension of its field of fractions, B a ring containing A and contained in \(\overline{\mathbf{A}}\). We set \(\mathbf{Y} = \operatorname{Spec}(\mathbf{B})\) and \(\mathbf{X} = \operatorname{Spec}(\mathbf{A})\) . We propose to prove that Y is a noetherian space. Let \((f_n)_{n \in \mathbb{N}}\) be a sequence of elements of B. For \(n \in \mathbb{N}\) , let \(\mathbf{U}_n = \operatorname{Spec}(\mathbf{B}_{f_n})\) be the open subset of Y defined by \(f_n\) . Put \(\mathbf{U} = \bigcup_{n \in \mathbb{N}} \mathbf{U}_n\) , \(\mathbf{F}_n = \mathbf{U} - \bigcup_{0 \leqslant m \leqslant n} \mathbf{U}_m\) . We denote by \(\overline{\mathbf{F}}_n\) the closure of \(\mathbf{F}_n\) in Y, \(\mathbf{G}_n\) the image of \(\mathbf{F}_n\)
\end{enumerate}

in X, \(\overline{\mathbf{G}}_n\) the closure of \(\mathbf{G}_n\) in X, \(\mathbf{B}_n\) the reduced quotient ring of B such that \(\operatorname{Spec}(\mathbf{B}_n) = \overline{\mathbf{F}}_n\) , \(\overline{f}_n\) the image of \(f_n\) in \(\mathbf{B}_n\) .

\begin{enumerate}
\def\labelenumi{\alph{enumi})}
\item
  We have \(\mathbf{F}_n = \mathbf{V}(\overline{f}_n\mathbf{B}_n)\cap \mathbf{U}\) for every \(n\in \mathbf{N}\) .
\item
  Suppose that for some integer \(n \in N\) , \(\overline{F}_{n}\) has only finitely many irreducible components. Then \(V(\overline{f}_{n}B_{n})\) has only finitely many irreducible components. (Apply the result of exerc. 15, also using exerc. 11, a).)
\item
  For every \(n\in\mathbf{N}\) , \(\mathbf{F}_{n}\) has only finitely many irreducible components.
\item
  If \(\mathbf{A}_n\) is the reduced quotient ring of A such that \(\operatorname{Spec}(\mathbf{A}_n) = \overline{\mathbf{G}}_n\) , then the minimal prime ideals of \(\mathbf{A}_n\) belong to \(\mathbf{G}_n\) .
\item
  Let \(n \in \mathbf{N}\) and \(\mathfrak{x}\) be such a minimal prime ideal of \(\mathbf{A}_n\) . Then there exists an integer \(n' > n\) such that \(\mathbf{F}_{n'}\) contains no point of \(\mathbf{Y}\) above \(\mathfrak{x}\) , and \(\mathfrak{x}\) does not belong to \(\overline{\mathbf{G}}_{n'}\) . (Use exerc. 11, a).)
\item
  For n sufficiently large, \(F_{n}\) is empty.
\end{enumerate}

\begin{enumerate}
\def\labelenumi{\arabic{enumi})}
\setcounter{enumi}{16}
\tightlist
\item
  Let A be a Noetherian integral domain and \(a \neq 0\) an element of the radical of A. We make the following hypotheses:
\end{enumerate}

\begin{enumerate}
\def\labelenumi{(\roman{enumi})}
\item
  \(\operatorname{Ass}_{\mathbf{A}}(\mathbf{A} / a\mathbf{A})\) contains a single minimal element \(p\) ;
\item
  \(a\mathbf{A}_{\mathfrak{p}} = \mathfrak{p}\mathbf{A}_{\mathfrak{p}}\)
\item
  A/p is integrally closed;
\item
  the integral closure \(\overline{\mathbf{A}}\) of A is a finite A-algebra;
\item
  If \(\mathfrak{q}\) is a prime ideal of height 1 in \(\overline{\mathbf{A}}\), then \(\mathfrak{q} \cap A\) is a prime ideal of height 1 in A.
\end{enumerate}

Prove that we have \(\overline{\mathbf{A}} = \mathbf{A}\) (A is integrally closed) and \(a\mathbf{A} = \mathfrak{p}\) (hence A/aA is integrally closed). To do this, establish successively the following assertions:

\begin{enumerate}
\def\labelenumi{\alph{enumi})}
\item
  \(A_{p}\) and \(\overline{A}_{p}\) are isomorphic and \(\overline{p} = p\overline{A}\) is the unique prime ideal of \(\overline{A}\) above p.
\item
  Additionally, \(\overline{p}\) is the unique minimal element of \(\operatorname{Ass}_{\mathbf{A}}(\overline{\mathbf{A}}/a\overline{\mathbf{A}})\) .
\item
  In fact, \(\overline{p}\) is the only element of \(\operatorname{Ass}_{\mathbf{A}}(\overline{\mathbf{A}}/a\overline{\mathbf{A}})\) .
\item
  \(\overline{\mathbf{A}} / a\overline{\mathbf{A}}\) is an integral domain and \(a\overline{\mathbf{A}} = \overline{\mathfrak{p}}\) .
\item
  A/p and \(\overline{A}/\overline{p}\) are isomorphic, and the canonical homomorphism from A/aA into \(\overline{A}/a\overline{A}\) is surjective.
\end{enumerate}

\begin{enumerate}
\def\labelenumi{\arabic{enumi})}
\setcounter{enumi}{17}
\tightlist
\item
  Let A be a Noetherian integral domain, \(x\) a nonzero element of A. We assume that A is separated and complete for the topology \(xA\) -adic and that for every prime ideal \(p \in \mathrm{Ass}_{\mathbf{A}}(\mathbf{A} / x\mathbf{A})\) , the ring \(A/p\) is Japanese. Prove that A is Japanese. (Let \(\overline{\mathbf{A}}\) be the integral closure of A in a finite extension of its field of fractions.
\end{enumerate}

\begin{enumerate}
\def\labelenumi{\alph{enumi})}
\item
  Since \(\overline{\mathbf{A}}\) is a Krull ring (p. 81, Exercise 14), let \(\mathfrak{P}_1, \ldots, \mathfrak{P}_n\) be the prime ideals of height 1 of \(\overline{\mathbf{A}}\) containing \(x\) .
\item
  Let \(\mathfrak{p}_i = \mathfrak{P}_i \cap \mathbf{A}\) . Then one has \(\mathfrak{p}_i \in \operatorname{Ass}_{\mathbf{A}}(\mathbf{A} / x\mathbf{A})\) and \(\kappa(\mathfrak{P}_i)\) is a finite extension of \(\kappa(\mathfrak{p}_i)\) (Use Exercises 11 and 13, p.~80.)
\item
  \(\overline{\mathbf{A}} / \mathfrak{P}_i\) is a finitely generated \(\mathbf{A} / \mathfrak{p}_i\)-module.
\item
  \(\overline{\mathbf{A}} / x\overline{\mathbf{A}}\) is a finitely generated A/xA-module.
\item
  \(\overline{\mathbf{A}}\) is separated for the topology \(x\overline{\mathbf{A}}\) -adic.
\end{enumerate}

Conclude as in the proof of Th. 1 of § 4, No. 2.

\begin{enumerate}
\def\labelenumi{\arabic{enumi})}
\setcounter{enumi}{18}
\tightlist
\item
  Let A be a Noetherian ring, and let a be a nonzero ideal of A. Suppose that A is separated and complete for the a-adic topology and that A/a is a Nagata ring. Prove that A is a Nagata ring. One may reason as follows:
\end{enumerate}

\begin{enumerate}
\def\labelenumi{\alph{enumi})}
\item
  Assume that A is an integral domain and that for every nonzero prime ideal p of A, the ring A/p is Japanese. Let \(\overline{\mathbf{A}}\) be the integral closure of A in a finite extension of its field of fractions. Then, for every nonzero prime ideal \(\mathfrak{P}\) of \(\overline{\mathbf{A}}\) , the ring \(\overline{\mathbf{A}}/\mathfrak{P}\) is a finite algebra over the ring \(A/(\mathfrak{P} \cap A)\). (Use exerc. 11, a), p. 80.)
\item
  With the same hypotheses and notations as in a), prove that \(\overline{A}\) is Noetherian. (Use Exercises 7, p. 79 and 14, p. 81.) Then, reasoning as in the previous exercise, prove that \(\overline{A}/a\overline{A}\) is a finitely generated A-module, that \(\overline{A}\) is separated for the topology \(a\overline{A}\) and that \(\overline{A}\) is a finitely generated A-module.
\end{enumerate}

\begin{enumerate}
\def\labelenumi{\arabic{enumi})}
\setcounter{enumi}{19}
\item
  Let A be a Noetherian ring, and let a be an ideal of A distinct from A, \(\hat{A}\) the separated completion of A for the a-adic topology. If A/a is a Nagata ring, then \(\hat{A}\) is a Nagata ring. (Use Exerc. 19.)
\item
  \begin{enumerate}
  \def\labelenumii{\alph{enumii})}
  \tightlist
  \item
    Let A be an integral domain, K its field of fractions. If K has nonzero characteristic p and E is a finite purely inseparable extension of the field of rational functions \(\mathrm{K}(\mathrm{X})\), there exists a finite purely inseparable extension F of K and a power q of p such that \(\mathrm{F}(\mathrm{X}^{1/q})\) is an extension of E.
  \end{enumerate}
\end{enumerate}

\begin{enumerate}
\def\labelenumi{\alph{enumi})}
\setcounter{enumi}{1}
\item
  Let A be a Nagata ring that is integrally closed. Then the polynomial ring A{[}X{]} is Japanese.
\item
  Let A be a Nagata ring that is integrally closed. Then every polynomial ring \(A[X_1, \ldots, X_n]\) in a finite number of indeterminates is a Nagata ring.
\end{enumerate}

\begin{enumerate}
\def\labelenumi{\arabic{enumi})}
\setcounter{enumi}{21}
\item
  Let A be a Noetherian ring, and let a be an ideal of A distinct from A. Suppose that A is separated and complete for the a-adic topology and that A/a is a Nagata ring. Then every ring of restricted formal power series in a finite number of indeterminates is a Nagata ring. (Use III, § 4, exerc. 7, and exercs. 20 and 21 above.)
\item
  Let A be a reduced Noetherian semi-local ring, and let R be its total ring of fractions, \(\hat{A}\) the completion of A. For \(\hat{A}\) to be reduced, it is necessary and sufficient that \(R \otimes_{A} \hat{A}\) is reduced.
\item
  Let A be a semi-local Noetherian ring, \(\hat{\mathbf{A}}\) its completion, \(x\) a nonzero element of the radical of A. We assume that \(\mathrm{Ass}_{\mathbf{A}}(\mathbf{A}/x\mathbf{A})\) contains no embedded prime ideal and that, for every \(\mathfrak{p} \in \mathrm{Ass}_{\mathbf{A}}(\mathbf{A}/x\mathbf{A})\) , the ring \(\mathbf{A}_{\mathfrak{p}}\) is regular and the ring \(\kappa(\mathfrak{p}) \otimes_{\mathbf{A}} \hat{\mathbf{A}}\) is reduced. Then \(\hat{\mathbf{A}}\) is reduced. (Reason as in the proof of Theorem 3, (i) \(\Rightarrow\) (ii) of § 4, No. 4.)
\item
  Let A be a Noetherian integral domain, X the topological space \(\operatorname{Spec}(A)\), and Nor(X) the set of points \(\mathfrak{p}\) of X such that the local ring \(A_{\mathfrak{p}}\) is integrally closed. (If the integral closure of A is a finite A-algebra, then Nor(X) is open in X, cf. V, § 1, No. 5, Cor. 5.)
\end{enumerate}

Let \(f\) be a nonzero element of A such that the ring \(\mathbf{A}[f^{-1}]\) is integrally closed.

\begin{enumerate}
\def\labelenumi{\alph{enumi})}
\item
  If a prime ideal p of A does not contain f, then p belongs to Nor(X).
\item
  Let E be the set of prime ideals \(\mathfrak{p}\) of A, associated with A/fA, of height \textgreater{} 1 or else of height 1 and such that \(A_{\mathfrak{p}}\) is not regular. Then we have
\end{enumerate}

\[
\operatorname{Nor} (X) = X - \bigcup_ {\mathfrak {p} \in E} V (\mathfrak {p}).
\]

(Use Exercises 4 and 5, p.~79.)

\begin{enumerate}
\def\labelenumi{\alph{enumi})}
\setcounter{enumi}{2}
\tightlist
\item
  Nor(X) is open in X.
\end{enumerate}

\begin{enumerate}
\def\labelenumi{\arabic{enumi})}
\setcounter{enumi}{25}
\tightlist
\item
  Let A be a Noetherian integral domain, X the topological space \(\operatorname{Spec}(A)\), Nor(X) the subspace of X introduced in Exercise 25, \(\overline{A}\) the integral closure of A. We assume that there exists a nonzero element \(f\) of A such that the ring \(A[f^{-1}]\) is integrally closed, and that for every maximal ideal m of A, the ring \(\overline{A}_{m}\) is a finite A \(_{m}\)-algebra. Consider \(\overline{A}\) as a filtered increasing inductive limit of finite sub-A-algebras, \(\overline{A} = \varinjlim(A_{j})_{j\in J}\). Set \(X_{j} = \operatorname{Spec}(A_{j})\) and let \(G_{j}\) be the image in X of \(X_{j} - \operatorname{Nor}(X_{j})\).
\end{enumerate}

\begin{enumerate}
\def\labelenumi{\alph{enumi})}
\item
  \(\operatorname{Nor}(\mathbf{X}_j)\) is open in \(\mathbf{X}_j\) . (Use Exerc. 25.)
\item
  \(G_{j}\) is closed in X.
\item
  There exists an index \(j_0 \in J\) such that one has \(G_{j_0} = \bigcap G_j\) .
\end{enumerate}

\(d)\) Let \(\mathfrak{p} \in \operatorname{Spec}(\mathbf{A})\) . Then for \(j \in \mathrm{J}\) sufficiently large, \(\mathbf{G}_j\) does not contain \(\mathfrak{p}\) .

\begin{enumerate}
\def\labelenumi{\alph{enumi})}
\setcounter{enumi}{4}
\tightlist
\item
  \(\overline{\mathbf{A}}\) is a finite A-algebra.
\end{enumerate}

\begin{enumerate}
\def\labelenumi{\arabic{enumi})}
\setcounter{enumi}{26}
\tightlist
\item
  Let A be a Noetherian integrally closed local ring. Assume moreover that A is a Nagata ring. Let x be an element of the fraction field K of A, not belonging to A. Let B be the ring A{[}x{]} and let p be a maximal ideal of B containing x. Set I = xA ∩ A.
\end{enumerate}

\begin{enumerate}
\def\labelenumi{\alph{enumi})}
\item
  Let X be an indeterminate and let Q be the kernel of the homomorphism of A{[}X{]} in B that maps X to x. Then Q is generated by the polynomials of the form aX - b, where a and b are elements of A such that ax = b, and I = xA ∩ A is generated by the constant terms b of these polynomials. Moreover, the inclusion of A into B induces an isomorphism from A/I onto B/xB.
\item
  The ideal I of A is divisorial.
\item
  \(\operatorname{Ass}_{\mathbf{B}}(\mathbf{B}/x\mathbf{B})\) contains no embedded prime ideal.
\item
  Let q be a prime ideal of \(B_{p}\) associated with \(B_{p}/xB_{p}\) . Then \(q \cap A\) is associated with A/I and the rings \(A_{q \cap A}\) and \((\mathbf{B}_{\mathfrak{p}})_{q}\) are discrete valuation rings.
\item
  \(B_{p}/q\) is a Nagata ring. (Note that \(\mathrm{B}/(\mathrm{q} \cap \mathrm{B})\) is isomorphic to \(\mathrm{A}/(\mathrm{q} \cap \mathrm{A})\) .)
\item
  The completion of \(B_{p}/q\) is reduced.
\item
  The completion of \(B_{p}\) is reduced and the integral closure of \(B_{p}\) is finite over \(B_{p}\) . (Use Exerc. 24.)
\end{enumerate}

\begin{enumerate}
\def\labelenumi{\arabic{enumi})}
\setcounter{enumi}{27}
\tightlist
\item
  Let A be a Noetherian local integrally closed ring, K its field of fractions, x an element of K - A, B the ring A{[}x{]} and let p be a maximal ideal of B. Let a, b be nonzero elements of A such that bx = a.
\end{enumerate}

\begin{enumerate}
\def\labelenumi{\alph{enumi})}
\item
  The ring B \(\left[\frac{1}{b}\right]\) is integrally closed.
\item
  There exists a monic polynomial \(f\) with coefficients in A such that \(f(x)\in\mathfrak{p}\) .
\item
  Let E be the field obtained by adjoining the roots of \(f(X)\) to K, A' the integral closure of A in E, B' the ring A'{[}x{]}. Let \(\mathfrak{p}'\) be a maximal ideal of B' lying over p. Then the integral closure of \(B'_{\mathfrak{p}'}\) is finite over \(B'_{\mathfrak{p}'}\). (Use the preceding exercise.)
\item
  The integral closure of \(\mathbf{B}_{\mathfrak{p}'}^{\prime}\) is finite over \(\mathbf{B}_{\mathfrak{p}'}^{\prime}\) . (Use exerc. 26.)
\item
  The integral closure of \(B_{p}\) is finite over \(B_{p}\) .
\end{enumerate}

\begin{enumerate}
\def\labelenumi{\arabic{enumi})}
\setcounter{enumi}{28}
\tightlist
\item
  Let A be an integrally closed Nagata ring, and x an element of the field of fractions of A, not belonging to A. Let B be the ring A{[}x{]} and let a, b be two nonzero elements of A such that bx = a.
\end{enumerate}

\begin{enumerate}
\def\labelenumi{\alph{enumi})}
\item
  The ring \(\mathbf{B}\left[\frac{1}{b}\right]\) is integrally closed.
\item
  For every maximal ideal m of B, the integral closure of \(B_{m}\) is a finite \(B_{m}\)-algebra. (Use the preceding exercise.)
\item
  The integral closure of B is a finite B-algebra. (Use exerc. 26.)
\end{enumerate}

\begin{enumerate}
\def\labelenumi{\arabic{enumi})}
\setcounter{enumi}{29}
\item
  Let A be a Nagata ring. Then every algebra of finite type over A is a Nagata ring. (Use exerc. 21, p. 82 and 29.)
\item
  Let A be a Noetherian integral domain and \(\overline{A}\) the integral closure of A in a finite extension of its field of fractions. If A is of dimension at most 2, then \(\overline{A}\) is a Noetherian ring. (Use exerc. 7, p. 79, exerc. 14, p. 81 and the Krull-Akizuki theorem (cf. VII, § 2, no. 5).)
\item
  Let A be a local Noetherian integral domain and K its field of fractions. Suppose K has nonzero characteristic p. Suppose that A is a Nagata ring and denote by \(\hat{A}\) the completion of A. Let p be a minimal prime ideal of \(\hat{A}\) and let L be the field of fractions of \(\hat{A}/p\) . Let k be a weak residue field of A, \(k'\) the field of representatives of \(\hat{A}\) containing the image of k in \(\hat{A}\) (p. 78, Exercise 9(a).) Then K is separable over k if and only if L is separable over \(k'\) . (Note that L is separable over K since A is a Nagata ring. If K is separable over k, prove that every derivation of \(k'\) into L extends to a derivation of L into L.)
\end{enumerate}

\section*{Exercises for the Appendix}
\phantomsection
\addcontentsline{toc}{section}{Exercises for the Appendix}

\begin{enumerate}
\def\labelenumi{\arabic{enumi})}
\tightlist
\item
  Let I be a nonempty right-filtering preordered set, and let \((\mathbf{A}_{\alpha}, \varphi_{\beta\alpha})\) an inductive system of rings relative to I. For each index \(\alpha\) in I, one is given an ideal \(q_{\alpha}\) of \(A_{\alpha}\) . Suppose that one has \(\varphi_{\beta\alpha}(q_{\alpha}) A_{\beta} = q_{\beta}\) for \(\beta \geqslant \alpha\) . We denote by A the inductive limit of \(A_{\alpha}\) , and for \(\alpha \in I\) , denote by \(\varphi_{\alpha}: A_{\alpha} \to A\) the canonical homomorphism. We set
\end{enumerate}

\[
\mathfrak {q} = \varinjlim \mathfrak {q} _ {\alpha}.
\]

\begin{enumerate}
\def\labelenumi{\alph{enumi})}
\item
  If for every \(\alpha \in I\) , \(q_{\alpha}\) is a maximal ideal of \(A_{\alpha}\) , then \(q\) is a maximal ideal of \(A\) .
\item
  If for every \(\alpha \in I\) , \(\mathfrak{q}_{\alpha}\) is contained in the radical of \(A_{\alpha}\) , then \(\mathfrak{q}\) is contained in the radical of \(A\) .
\item
  If for every \(\alpha\in\mathbf{I}\) and every \(\beta\in\mathbf{I}\) , such that \(\beta\geqslant\alpha\) , the homomorphism \(\varphi_{\beta\alpha}\) is faithfully flat, then for all \(\alpha\in\mathbf{I}\) , \(\varphi_{\alpha}\) is faithfully flat.
\end{enumerate}

We shall henceforth assume that the hypothesis of c) is satisfied.

\(d)\) If for every \(\alpha\in I\) , \(\mathbf{A}_{\alpha}\) is separated for the topology \(\mathfrak{q}_{\alpha}\) -adic, then A is separated for the topology \(\mathfrak{q}\) -adic.

\begin{enumerate}
\def\labelenumi{\alph{enumi})}
\setcounter{enumi}{4}
\tightlist
\item
  If for every \(\alpha \in I\) , \(A_{\alpha}\) is Noetherian and if A/q is Noetherian, then \((1 + q)^{-1}A\) is Noetherian. (Reason as in the proof of Prop. 1 of the Appendix, additionally using the following result, which one will prove:
\end{enumerate}

Let B be a commutative ring, and let p be an ideal of B, \(\widehat{B}\) the separated completion of B for the p-adic topology, and \(i: B \to \widehat{B}\) the canonical map. Then there exists one and only one ring homomorphism \(j: (1 + p)^{-1}B \to \widehat{B}\) which extends \(i\). For every maximal ideal m of \((1 + p)^{-1}B\) , one has \(j(m)\widehat{B} \neq \widehat{B}\) .

\begin{enumerate}
\def\labelenumi{\arabic{enumi})}
\setcounter{enumi}{1}
\item
  Let A be a local ring, and B a blowing-up of A. If A is separated (for the topology defined by the maximal ideal \(m_{A}\)), B is separated (for the topology defined by the maximal ideal \(m_{B}\)).
\item
  Let \(k\) be an imperfect field of characteristic \(p > 0\) and \(\mathbf{R} = k[\mathrm{T}]_{(\mathrm{T})}\) . Let \(a\) be an element of \(k\) which is not a \(p\)-th power. Show that the ring \(\mathbf{A} = \mathbf{R}[\mathbf{X}] / (\mathbf{X}^p - a\mathbf{T}^p)\) is local, an integral domain, and has residue field \(k\) . Show that the ring \(\mathbf{B} = \mathbf{A}[\mathbf{Y}] / (\mathbf{Y}^p - a)\), which is a blowing-up of \(\mathbf{A}\), is not reduced.
\item
  Let A be a local ring equipped with a valuation \(v_{\mathbf{A}}\) . Let B be a blowing-up of A. Show that the valuation \(v_{\mathbf{A}}\) extends uniquely to a valuation \(v_{\mathbf{B}}\) of B, and that the value groups \(v_{\mathbf{A}}(\mathbf{A})\) and \(v_{\mathbf{B}}(\mathbf{B})\) coincide. If A is a discrete valuation ring with uniformizer \(\pi\), B is also a discrete valuation ring, one of whose uniformizers is the image of \(\pi\) in B.
\item
  Let A be a local ring with residue field \(\kappa(A)\) . Let B be a blow-up of A, with residue field \(\kappa(B)\) .
\end{enumerate}

\begin{enumerate}
\def\labelenumi{\alph{enumi})}
\item
  If \(\kappa(\mathbf{B}) = \kappa(\mathbf{A})\) , then \(\mathbf{A} = \mathbf{B}\) .
\item
  If \(\kappa(\mathbf{B})\) is algebraic over \(\kappa(\mathbf{A})\) , the A-algebra B is integral.
\item
  If \(\kappa(B)\) has finite degree over \(\kappa(A)\) , \(B\) is a free A-module whose rank is the degree of \(\kappa(B)\) over \(\kappa(A)\) .
\end{enumerate}

\begin{enumerate}
\def\labelenumi{\arabic{enumi})}
\setcounter{enumi}{5}
\tightlist
\item
  Let A be a local ring with residue field \(k\) . Let I be a set of indices. Show that the ring \(\mathbf{A}](\mathbf{X}_i)_{i\in \mathbf{I}}[\) (cf. Appendix, Example 2), which is a blowing-up of A, has as its residue field the purely transcendental extension \(k((\mathbf{X}_i)_{i\in \mathbf{I}})\) of \(k\) .
\end{enumerate}

\specialchapter{Index of notations}

\indexgroup{Chapter VIII}

\(\kappa(p), m_A, \kappa_A, SM: VIII, p. 1.\)

dim kr(X), dim(X): VIII, p.~2.

\[
\dim_ {x} (X): \text { VIII }, p. 2.
\]

codim(Y,X) : VIII, p.~4.

\[
\dim (\mathrm{A}): \text { VIII }, \text { p. } 6.
\]

\[
\dim_ {\mathfrak {p}} (\mathrm{A}): \text { VIII }, \text { p. } 6.
\]

ht(a): VIII, p.~8.

\(\dim_{\mathbf{A}}(\mathbf{M}), \dim(\mathbf{M}) : \mathrm{VIII}, \mathrm{p.} 10.\)

\[
[ \mathrm{Y} ], \bar {\mathrm{Z}} (\mathrm{A}): \text { VIII }, \text { p. } 11.
\]

\[
z (\mathbf {M}), \mathbf {Z} _ {\leqslant d}, \mathbf {Z} _ {d}, \mathbf {Z} ^ {\geqslant d}, \mathbf {Z} ^ {d}: \text { VIII }, \text { p. } 11.
\]

\[
\mathbb {C}, \mathbb {C} _ {\leqslant d}, \mathbb {C} ^ {\geqslant d}, \zeta_ {d}, \zeta^ {d}: \text { VIII }, \text { p. } 12.
\]

\[
(\mathrm{PM}) (\text { condition }): \text { VIII }, \text { p. } 13.
\]

\[
\mathrm{A} ((\mathrm{T})), \left[ \begin{array}{c} n \\ p \end{array} \right]: \text {VIII, p. 38.}
\]

\(\mathrm{F} \leqslant \mathrm{G}\) (for F and G in \(\mathbf{Z}((\mathrm{T}))\) ): VIII, p.~39.

\[
\mathrm{P} _ {\mathrm{M}}, \mathrm{Q} _ {\mathrm{M}}: \mathrm{VIII}, \text {p.} 40.
\]

\[
c _ {\mathbf {M}}: \text { VIII }, \text { p. } 41.
\]

\[
\tilde {\mathbf {G}} ^ {(r)}: \text { VIII }, \text { p. } 43.
\]

\[
\mathrm{H} _ {\mathbf {M}, \mathbf {F}}: \text { VIII }, \text { p. } 44.
\]

\[
\mathrm{H} _ {\mathrm{M}, \mathrm{q}}: \text { VIII }, \text { p. } 44.
\]

\[
\mathrm{d} _ {\mathfrak {q}} (\dot {\mathbf {M}}), \mathrm{e} _ {\mathfrak {q}} (\mathbf {M}): \text { VIII }, \text { p. } 45.
\]

\[
\mathrm{H} _ {\geqslant n}: \text { VIII }, \text { p. } 62.
\]

dim gr(H), ht gr(p): VIII, p.~63.

\[
\mathfrak {a} ^ {g r}: \text { VIII }, \text { p. } 63.
\]

\[
\mathbf {M} _ {\geqslant n}: \text { VIII }, \text { p. } 65.
\]

\[
\mathrm{H} _ {+}: \text {VIII}, \text {p.} 70.
\]

\(e_{\mathfrak{q}}^{\mathbf{A}}(\mathbf{M}), e^{\mathbf{A}}(\mathbf{M}), e(\mathbf{M}) : \mathrm{VIII}, \mathrm{p.72}\) .

\(R_{M}\) VIII, p. 77.

Specmax(A): VIII, p.~81, exerc. 2.

dev(E): VIII, p.~81, exerc. 4.

Kdev(M): VIII, p.~81, exerc. 5.

\(\dim_{\mathbf{v}}(\mathbf{A}) : \mathbf{VIII}, \text{ p. } 85, \text{ exerc. } 16.\)

\(\partial_l(n), \partial_l^*(n), S_H(T): VIII, p. 90, exerc. 7.\)

Reg(A), Sing(A): VIII, p.~96, exerc. 16.

Lis(A): VIII, p.~98, exerc. 25.

\(\overline{\mathfrak{a}}:\mathbf{V}\mathbf{I}\mathbf{I}\mathbf{I}\mathbf{,}\mathbf{p.}105,\) Exercise 9.

Chapter IX

\[
\Phi_ {n}: \mathrm{IX}, \text { p. } 1.
\]

\[
f _ {\mathrm{A}}, v _ {\mathrm{A}}, f, v: \mathrm{IX}, \text { p. } 2.
\]

\[
\Phi_ {A}, \Phi : \mathrm{IX}, p. 3.
\]

\[
\mathrm{S} _ {n}, \mathrm{P} _ {n}, \mathrm{I} _ {n}, \mathrm{F} _ {n}: \mathrm{IX}, \text { p. } 4.
\]

\[
\mathrm{S} _ {\mathrm{A}}, \mathrm{P} _ {\mathrm{A}}, \mathrm{I} _ {\mathrm{A}}, \mathrm{W} (\mathrm{A}), \mathrm{W} (\rho), \rho^ {\mathrm{N}}: \mathrm{IX}, \text { p. } 6.
\]

\[
\mathrm{F} _ {\mathrm{A}}, \mathrm{V} _ {\mathrm{A}}, \mathrm{F}, \mathrm{V}: \mathrm{IX}, \text {p. 8.}
\]

\[
\mathcal {C}, \mathrm{V} _ {n} (\mathrm{A}), \tau_ {\mathrm{A}}, \tau : \mathrm{IX}, \text { p. } 11.
\]

\[
\mathrm{W} _ {n} (\mathrm{A}): \mathrm{IX}, \text { p. } 12.
\]

\[
\mathrm{W} _ {n} (\rho), \pi_ {n}, \pi_ {n, m}: \mathrm{IX}, \text { p. } 12.
\]

\[
\mathrm{V} _ {m} ^ {n}, \mathrm{F} _ {m} ^ {n}: \mathrm{IX}, \text { p. } 14.
\]

\[
l (\mathbf {C}): \mathrm{IX}, \text { p. } 18.
\]

\[
\mathrm{A} ] \mathrm{X} [: \mathrm{IX}, \mathrm{p}. 37.
\]

\(\mathrm{A}]\left(\mathrm{X}_{i}\right)_{i\in \mathbf{I}}[: \mathrm{IX},\mathrm{p.}39.\)

\(s_{\mathrm{A}}:\mathrm{IX},\mathrm{p.}44,\) Exercise 15.

\(\rho,\mathrm{K}(\rho^{-1}(\mathrm{A})),[\sigma,a>:\mathrm{IX},\mathrm{p.~45,~exerc.~19}.\)

\[
\mathfrak {D} _ {\mathrm{A}}: \mathrm{IX}, \text { p. } 47, \text { exerc. } 22.
\]

\(\mathbf{CW}^u (\mathbf{A}):\mathbf{IX},\) p.~48, exerc. 23.

CW(A): IX, p.~48, exerc. 24.

J, \(\Phi, f_n, v_n, \Phi_n: \mathrm{IX}\) , p.~50, exerc. 28.

\(\mathbf{U}(\mathbf{A}), \mathbf{F}_n, \mathbf{V}_n, \mathbf{U}(\rho) : \mathbf{IX}, \text{p. 52, exerc. 35 and 36.}\)

\(\tau_{\mathrm{A}}:\mathrm{IX},\mathrm{p.~53,~exerc.~37.}\)

\[
\mathrm{U} _ {S} (\mathrm{A}), \mathrm{U} _ {n} (\mathrm{A}); \mathrm{U} _ {n \infty} (\mathrm{A}): \mathrm{IX}, \text { p. } 53, \text { exerc. } 39.
\]

\(\mu^{\mathbf{A}}, \mathbf{F}^{\mathbf{A}}, f^{\mathbf{A}}: \mathrm{IX}, \mathrm{p.~55,~exerc.~41}\) .

\(\Lambda (\mathbf{A}),\mathbf{L}:\mathbf{IX},\mathbf{p}.56,\) exerc. 42.

\(\mathrm{Rep}_{\mathbf{C}}(\mathbf{A}), \mathbf{R}_{\mathbf{C}}(\mathbf{\hat{A}}) : \mathbf{IX}, \text{p. 64, exerc. 54.}\)

\(\operatorname{Rep}_{\mathrm{C}}(\mathbf{G}), \mathbf{R}_{\mathrm{C}}(\mathbf{G}) : \mathrm{IX}, \mathrm{p.~65,~exerc.~55}.\)

\(\delta (\mathbf{T}),\mathrm{E}_{\sigma}(a,\mathbf{T}):\mathrm{IX},\mathrm{p.~68,~exerc.~58.}\)

\(A_{red}\) IX, p.~80, exerc. 9.

Nor(X): IX, p.~83, exerc. 25.

\specialchapter{Index of terminology}

Regular graded algebra: VIII, p.~71. Catenary ring: VIII, p.~9. Nagata ring: IX, p.~34. Ring of Witt vectors: IX, p.~7. Universal Witt ring: IX, p.~7. Equidimensional ring: VIII, p.~82, exerc. 10. Japanese ring: IX, p.~30. Local ring of equal characteristics: IX, p.~28. Regular local ring: VIII, p.~52. Regular Noetherian ring: VIII, p.~94, exerc. 6. Normal ring: IX, p.~78, exerc. 3.

Catenary (ring): VIII, p.~9. Catenary (space): VIII, p.~5. Chain: VIII, p.~1. Codimension (of Y in X): VIII, p.~4. Ghost component: IX, p.~7 and p.~12. Field of representatives: IX, p.~29. Witt covectors: IX, p.~48, exerc. 24. Unipotent Witt covectors: IX, p.~48, exerc. 23. Cycle: VIII, p.~11.

Shift: IX, p.~8. Deviation (of an ordered set): VIII, p.~81, exerc. 4. (Krull) dimension of a ring: VIII, p.~6. Dimension of a ring at a prime ideal: VIII, p.~6. (Krull) dimension of a space: VIII, p.~2. Dimension of a module: VIII, p.~10. Valuative dimension: VIII, p.~85, exerc. 16.

Eisenstein (polynomial): VIII, p.~56. Equidimensional (ring): VIII, p.~82, exerc. 10. Staircase: VIII, p.~92, exerc. 8. Extremities of a chain: VIII, p.~1.

Secant family: VIII, p.~26. Integral closure of an ideal: VIII, p.~105, exerc. 9. Hilbert-Samuel function of a filtered module: VIII, p.~44. Macaulay function: VIII, p.~90, exerc. 7. Formally smooth (algebra): VIII, p.~98, exerc. 30 and IX, p.~76, exerc. 2.

Swelling of a local ring: IX, p.~38. Elementary swelling of a local ring: IX, p.~38.

Height of an ideal: VIII, p.~8. Hironaka (lemma of): VIII, p.~107, exerc. 22. Hypersurface: VIII, p.~29.

Prime ideal lying over another: VIII, p.~1. Complete intersection: VIII, p.~103, exerc. 3.

Singular locus: VIII, p.~96, exerc. 16. Smooth (algebra): IX, p.~76, exerc. 2. Smooth (prime ideal of an algebra): VIII, p.~98, exerc. 23. Length of a chain: VIII, p.~1. Length of a p-ring: IX, p.~18.

λ-ring: IX, p.~59, exerc. 45. λ-ideal: IX, p.~63, exerc. 50. λ-morphism: IX, p.~58, exerc. 45. λ-rank: IX, p.~59, exerc. 46.

Macaulay (function of): VIII, p.~90, exerc. 7. Macaulay (theorem of): VIII, p.~93, exerc. 11. Multiplicity (of a module with respect to an ideal): VIII, p.~72.

Normal (ring): IX, p.~78, exerc. 3.

p-ring: IX, p.~17. Poincaré (series of): VIII, p.~40. Eisenstein polynomial: VIII, p.~56. Polynomial at infinity (map): VIII, p.~88, exerc. 2. Witt polynomials: IX, p.~1. Pre-λ-ring: IX, p.~58, exerc. 45.

Regular (local ring): VIII, p.~52. Regular (Noetherian ring): VIII, p.~94, exerc. 6. Regular (graded algebra): VIII, p.~71. Multiplicative representatives: IX, p.~25.

Saturated (chain): VIII, p.~1. Secant (family, sequence): VIII, p.~26. Poincaré series (of a graded module): VIII, p.~40. Hilbert-Samuel series (of a filtered module): VIII, p.~44. Singular (locus): VIII, p.~96, exerc. 16. Cohen subring (of a local ring): IX, p.~20. Secant sequence: VIII, p.~26. Superficial (element): VIII, p.~79. Superficial of order δ (element): VIII, p.~79. System of coordinates: VIII, p.~52. Graded system of coordinates: VIII, p.~71.

Chevalley theorem: VIII, p.~101, exerc. 10. Nagata theorem: IX, p.~32.

Witt vector: IX, p.~7. Witt vector of finite length: IX, p.~12.
